\chapter*{Preface}
\addcontentsline{toc}{chapter}{Preface}
\markboth{\textsc{Preface}}{}

There is a moment, somewhere around the middle of this book, when a small
plastic card and a three-pound circuit board prove to each other that they both
know the same sixteen-byte number --- without either of them ever transmitting
it, and while anyone within a few centimetres could be recording every bit that
passes between them.

That moment is the reason this book exists. It is the same handshake that
happens when you tap your bank card, open your office door, or board a train.
It is not magic and it is not beyond you. It is about two hundred lines of code,
and by the end of Chapter~\ref{ch:auth} you will have written all of them.

\section*{What you will build}

A coffee kiosk that takes contactless payments. Touch a drink on a
touchscreen, hold a card to the reader, and:

\begin{enumerate}
  \item the reader wakes the card out of the radio field and reads its serial number;
  \item your code and the card authenticate each other with AES-128;
  \item your code reads the balance --- encrypted, and authenticated, so a
        forged reply is rejected rather than believed;
  \item it debits the price, commits the transaction, and the drink is dispensed.
\end{enumerate}

Plus a phone app that puts credit back on.

Every layer is yours. No NFC library, no crypto library beyond a raw AES block
cipher, no copied driver. You will write the SPI register encoding, the
anticollision tree walk, the ISO-DEP block framing, AES-CMAC from the RFC, the
session key derivation, the mutual authentication, and the secure messaging that
carries the money.

\section*{Who this is for}

You know C. You have blinked an LED, read a datasheet, and written
\code{uint8\_t buf[16]}. You have probably not written much Java, and you may
never have seen an interface, an abstract class, or a line of cryptography.

That is exactly the intended reader. This book teaches object-oriented design
and applied cryptography \emph{as tools for a job you are actually doing}, in the
order the job demands them.

\begin{concept}{No concept before the problem that needs it}
You will not read ``an interface is a contract between a class and its
callers'' as an abstract definition. You will hit a concrete wall --- \emph{I
cannot test my driver without a Raspberry Pi physically plugged in} --- and the
interface will be the tool that gets you over it.

Concepts learned that way stick, because they are attached to a memory of being
stuck. Concepts learned as definitions evaporate before the exam.
\end{concept}

Each chapter names the one design idea it introduces:

\begin{center}\small
\begin{tabular}{@{}ll@{}}
\toprule
\thead{Chapter} & \thead{Idea, and the problem that forces it} \\
\midrule
1  & Interfaces --- testing a driver with no hardware \\
2  & Encapsulation --- forty copies of a bit-twiddle \\
3  & Abstract classes --- one protocol, many reader chips \\
4  & Polymorphism --- three frame types, one call site \\
5  & Value objects --- three results that belong together \\
6  & Layering --- a seam that makes crypto testable \\
7  & Immutability --- an object that cannot be half-built \\
8  & Enums over booleans --- \code{derive(key, a, b, true)} means nothing \\
9  & Injecting randomness --- testing a function that is random \\
10 & Growing a class --- state and its operations belong together \\
11 & The Adapter pattern --- two fixed interfaces, made to fit \\
12 & Porting a design --- discovering what was incidental \\
\bottomrule
\end{tabular}
\end{center}

\section*{How to use this book}
\label{sec:howtouse}

\textbf{Read the chapter, close the book, write the code.} Every chapter gives
you five things:

\begin{enumerate}
  \item \textbf{The problem.} What does not work yet, concretely.
  \item \textbf{The idea.} The design, with an analogy from outside programming.
  \item \textbf{The boilerplate.} A skeleton with \code{TODO} comments.
  \item \textbf{The verification gate.} A test with known-good values.
        \emph{Do not continue until it passes.}
  \item \textbf{The traps.} The specific bugs that will bite, with their symptoms.
\end{enumerate}

\subsection*{Why the listings have holes in them}

Because reading code teaches you very little and writing it teaches you almost
everything.

The skeletons hand you the parts that are pure lookup --- register addresses,
class declarations, import lists, byte layouts --- and withhold the parts that
are actual reasoning: the algorithm, the control flow, the order of operations.
You spend your thinking where thinking pays.

If you get genuinely stuck, a complete implementation of every file exists in
the companion repository. Treat it as the back of a textbook: consulting it ends
the exercise, so spend twenty minutes stuck first. \emph{Being stuck is the
learning.} The answer merely ends it.

\subsection*{The verification gate, and why it is not optional}

Cryptographic code has a property that no code you have written before has:
\textbf{wrong output is indistinguishable from right output.}

A broken checksum gives you a visible mismatch. A broken CMAC gives you sixteen
bytes of perfect-looking randomness --- exactly like a correct CMAC. You cannot
eyeball it. You cannot infer it from a failed transaction. And if you build three
more layers before testing, you will have no way to know which one is lying.

On real hardware it is worse: a card may count consecutive authentication
failures and, after enough of them, lock a key permanently.

So every chapter in Part~IV ends with \textbf{known-answer test vectors} ---
inputs with independently computed correct outputs. Reproduce them and your code
is right. Fail them and it is wrong, and you fix it \emph{now}.

The vectors in this book were generated with Python's \code{cryptography}
library, deliberately \emph{not} with the Java you are writing. Otherwise they
would only prove your code agrees with itself.

\section*{What you need}

\begin{center}\small
\begin{tabular}{@{}L{0.34\linewidth}L{0.46\linewidth}r@{}}
\toprule
\thead{Item} & \thead{Notes} & \thead{Approx.} \\
\midrule
Raspberry Pi Zero 2 W & Any Pi works; the Zero 2 W is the tight fit that teaches you to optimise & \pounds18 \\
MFRC522 reader module & The ubiquitous breakout board & \pounds3 \\
\textbf{MIFARE DESFire EV2/EV3} & \textbf{Not Classic, not NTAG. See below.} & \pounds3 ea \\
7-inch touchscreen \emph{(optional)} & 1024$\times$600, HDMI + USB touch & \pounds45 \\
Jumper wires, microSD, 5\,V PSU & & \\
An Android phone with NFC & For Part~VI & \\
\bottomrule
\end{tabular}
\end{center}

\begin{trap}{The card matters more than anything else on that list}
It must say \textbf{DESFire EV2} or \textbf{EV3}. Not MIFARE Classic (broken
crypto, entirely different protocol), not NTAG or Ultralight (no crypto at all,
no value files), not ``13.56\,MHz RFID tag'' from a bag of a hundred.

If a listing says ``compatible'' or does not name the chip, it is not one.
\textbf{Buy at least two} --- one to develop against and one spare, because it
is genuinely possible to destroy a card by changing a key and forgetting it.

You can confirm what you have in Chapter~\ref{ch:apdu}: a real DESFire answers
\code{GetVersion} with a hardware type byte of \hex{0x33} for EV3.
\end{trap}

\section*{Conventions}

\begin{itemize}
  \item Java 21 and Maven, package \code{com.midnightbrewer}. Keeping the same
        package as the reference repository means you can diff against it.
  \item Hexadecimal is lowercase and unspaced in code, spaced in prose:
        \hex{91 00}.
  \item Coloured rules in the margin mark the four kinds of aside:
        \textcolor{brewerblue}{concepts}, \textcolor{breweramber}{analogies},
        \textcolor{brewerred}{traps}, and \textcolor{brewergreen}{verification
        gates}.
  \item Every chapter ends with a suggested commit message, so your history
        becomes a record of the build.
\end{itemize}

\vspace{1em}
\begin{flushright}
\itshape One layer at a time.
\end{flushright}
